\section{Married Priests in the Latin Church}

The Latin Church ordains married men to the presbyterate. She does it rarely, by a reserved act, and under a rule that names itself as a departure from the general norm---but she does it, and the legal instruments are public. Anyone who describes the Latin discipline as an absolute rule has described something the Church herself does not claim.

\subsection{The 1967 opening}

The provision was made in the same encyclical that reaffirmed the law. \emph{Sacerdotalis caelibatus} n.~42 states two things at once, and its Latin marks the pairing with \emph{hinc\ldots illinc}, on the one hand and on the other:

\begin{studyframe}
\latin{Quapropter ex primaria illa norma, quam ad regimen catholicae Ecclesiae spectantem supra commemoravimus, duo hoc loco esse statuenda censemus: hinc legem, quae eos, qui in sacros ordines ascribantur, caelibatum libere ac perpetuo eligere iubet, firmam in suo iure manere; illinc licere peculiares perspicere condiciones sacrorum administrorum, qui iam matrimonio coniuncti sive in ecclesiis sive in christianis communitatibus a catholica communione adhuc distinctis vivunt, si, plena huiusmodi communione frui sacroque postea ministerio fungi exoptantes, ad sacerdotalia officia vocentur; ea tamen ratione quae statutae iam disciplinae sacri caelibatus, a clero servandi, non obsit.}

\emph{Official English:} ``\ldots while on the one hand, the law requiring a freely chosen and perpetual celibacy of those who are admitted to Holy Orders remains unchanged, on the other hand, a study may be allowed of the particular circumstances of married sacred ministers of Churches or other Christian communities separated from the Catholic communion, and of the possibility of admitting to priestly functions those who desire to adhere to the fullness of this communion and to continue to exercise the sacred ministry. The circumstances must be such, however, as not to prejudice the existing discipline regarding celibacy.''
\end{studyframe}

Two witnesses were used for that Latin and they agree word for word: the encyclical as printed in \emph{Acta Apostolicae Sedis} 59 (1967) 674, and the Holy See's Latin web delivery. The English is the Holy See's own English delivery.\endnote{\emph{Sacerdotalis caelibatus} (24 June 1967) n.~42. Latin: \emph{Acta Apostolicae Sedis} 59 (1967) 674, read at the text layer of the Holy See's archival PDF of the volume, fetched and hashed 2026-07-25; that PDF's page rasters could not be rendered by the tools available here, so the \emph{Acta} reading is a text-layer reading and not a page-image reading, and this article says so. Corroborating witness: the Holy See's Latin web delivery of the encyclical, fetched, hashed, and read the same day, which agrees verbatim at every word of n.~42 except that it prints \latin{posset conferr} for \latin{posset conferri} in the following sentence---a defect of that delivery, recorded and not repaired. English from the Holy See's English web delivery, same date. The encyclical's authority is that of an encyclical letter: authoritative non-definitive papal teaching, which in n.~42 also exercises the governing power it describes.} The paragraph continues by noting that the Church's authority does not hesitate to act in this matter, ``as can be seen from the recent Ecumenical Council, which foresaw the possibility of conferring the sacred order of the diaconate on men of mature age who are married.''

That last clause matters for this article's spine. The Council's restoration of a married diaconate and the encyclical's opening for married convert ministers are presented by Paul~VI as instances of a single power: the Church's power over her own determination about who receives orders. Not a power over continence, and not a power over the impediment of canon 1087, neither of which n.~42 touches.

\subsection{\emph{Anglicanorum coetibus} and its complementary norms}

Forty-two years later the provision became structural. Benedict~XVI's apostolic constitution \emph{Anglicanorum coetibus} of 4 November 2009 erected personal ordinariates for groups of Anglicans entering full communion, and its article~VI governs their clergy.

Article VI §1 provides that those who ministered as Anglican deacons, priests, or bishops, and who fulfil the requisites of canon law and are not impeded by irregularities or other impediments, may be accepted by the Ordinary as candidates for Holy Orders in the Catholic Church; ``in the case of married ministers, the norms established in the Encyclical Letter of Pope Paul~VI \emph{Sacerdotalis coelibatus}, n.~42 and in the Statement \emph{In June} are to be observed''; and unmarried ministers ``must submit to the norm of clerical celibacy of CIC can.~277, §1.''

Article VI §2 states the rule and the exception in one sentence, and the Latin is precise where the English is merely clear: \latin{Ordinarius, omnino disciplinae in Ecclesia Latina circa caelibatum clericalem satisfaciens, pro regula ad presbyteralem ordinem dumtaxat viros admittet caelibes. A Romano Pontifice expetere poterit, can.~277, §~1 derogando, ut singulis in casibus, ad Ordinem Sacrum presbyteratus admittantur etiam coniugati viri, persedulo cautis tamen obiectivis criteriis ab Apostolica Sede comprobandis}---the Ordinary, in full observance of the discipline of celibate clergy in the Latin Church, ``as a rule (\latin{pro regula}) will admit only celibate men to the order of presbyter,'' and may petition the Roman Pontiff, ``as a derogation from can.~277, §1,'' for the admission of married men case by case, according to objective criteria approved by the Holy See.\endnote{\emph{Anglicanorum coetibus} art.~VI, Latin and English quoted from the Holy See's Latin and English web deliveries of the constitution, both fetched, hashed, and read 2026-07-25. The constitution's own note~15 cites \emph{AAS} 59 (1967) 674 for \emph{Sacerdotalis caelibatus} n.~42---the page verified above---and its note~16 identifies the ``Statement \emph{In June}'' as a Congregation for the Doctrine of the Faith statement of 1 April 1981, in \emph{Enchiridion Vaticanum} 7, 1213. That statement was not examined for this article and nothing here rests on its content; the discrepancy between the body's incipit-style name and the note's date is reported as the document prints it. The constitution's spelling \emph{Sacerdotalis coelibatus} in the English text, against \emph{caelibatus} in the Latin, is likewise as printed.}

The word \latin{derogando} is the technical point of the whole provision. A derogation is the partial removal of a law for a case; it is not a dispensation from an obligation already incurred, and it is not an amendment of the canon. Canon 277 §1 remains the law; the Roman Pontiff removes its application to a named man before he becomes a cleric. And what is derogated from is canon 277 §1 \emph{entire}---continence and celibacy together---which is the only coherent reading, since a married ordinand plainly is not bound to either half.

The Complementary Norms issued by the Congregation for the Doctrine of the Faith on the same day add three operative limits at their own article~6. The Ordinary needs the consent of his Governing Council to admit candidates to orders. The objective criteria for a married-clergy petition are ``determined by the Ordinary in consultation with the local Episcopal Conference and must be approved by the Holy See.'' And two categories are excluded outright: those previously ordained in the Catholic Church who subsequently became Anglicans may not exercise sacred ministry in an ordinariate, and ``Anglican clergy who are in irregular marriage situations may not be accepted for Holy Orders in the Ordinariate.''\endnote{Congregation for the Doctrine of the Faith, \emph{Complementary Norms for the Apostolic Constitution Anglicanorum coetibus}, art.~6 §§1--2, quoted from the Holy See's English delivery, fetched, hashed, and read 2026-07-25. The delivery carries a footnote recording that art.~5 §2 ``was added to the text of the Complementary Norms according to a decision of the Ordinary Session held on 29 May 2013, approved by Pope Francis on 31 May 2013,'' so the page delivers the amended text; the amendment does not touch art.~6. The Norms are an act of a dicastery under the constitution, subordinate to it, and---like the constitution---particular law for the ordinariates, not universal law.}

\subsection{What this does and does not establish}

It establishes that a married Latin presbyterate exists as a lawful, published, ordinary-course provision with named criteria and a named authority. Anyone arguing about the Latin discipline has to argue about the discipline that actually exists, which already contains this.

It does not establish that the Latin Church has changed her general norm, and the texts are careful to say so twice: \emph{Sacerdotalis caelibatus} 42's \latin{firmam in suo iure manere} and \emph{Anglicanorum coetibus} VI §2's \latin{pro regula}. Nor does it establish anything at all about ordaining married men who are not former ministers of another ecclesial communion; every instrument here is addressed to that specific ecumenical situation, and a norm addressed to a situation reaches that situation. Nor, finally, does it touch the impediment of canon 1087: an ordinariate presbyter widowed after ordination stands exactly where any other Latin cleric stands.

There is one further category the universal law recognizes but this article has not examined: the earlier ``Pastoral Provision'' arrangements for former Anglican clergy in the United States from 1980, and comparable case-by-case admissions of married former Lutheran and other ministers. They are named in \emph{Anglicanorum coetibus}'s own footnote and are real, and no claim is made here about their terms.
